\section{Proofs}
\label{app:proofs}

Throughout, ``result'' means the observable part of
\(\llbracket P \rrbracket(\sigma)\), and \(\mathit{Acc}_\theta(\sigma)\) is
the set of observable results that \(\mathit{Post}_\theta\) accepts on
\(\sigma\).

\paragraph{Proposition~\ref{prop:refine} (As-If Programs Remain Legal).}
If \(P \models \Phi_{\text{func}}\), \(P' \approx P\), and
\(\theta_{P'} = \theta_P\), then \(P' \models \Phi_{\text{func}}\).
\begin{proof}
Suppose \(P' \not\models \Phi_{\text{func}}\).
Then either \(\theta_{P'} \notin \mathcal{I}\), or some
\(\sigma \in \mathit{Pre}\) is a non-terminating, undefined, or unacceptable
run of \(P'\). The first case is impossible because
\(\theta_{P'} = \theta_P \in \mathcal{I}\). In the second case,
\(P' \approx P\) gives the same observable result for \(P\) on \(\sigma\),
so \(P\) also violates the contract, a contradiction.
\end{proof}

\paragraph{Proposition~\ref{prop:expansion} (Legal Set Expands Exactly on Observable Slack).}
Assume local patching, and that each \(\mathit{Post}_\theta\) with
\(\theta \in \mathcal{I}\) is realized by some program. Let
\(P \models \Phi_{\text{func}}\). Then
\(\mathrm{Sat}(\Phi_{\text{func}})\) contains a program that is not
observationally equivalent to \(P\) iff the contract has input, output, or
logical-parameter slack, counting only results some program produces.
Unexposed slack is not in that list: it need not change those results.
\begin{proof}
Assume also that all results below are produced by some program.
Let \(o_P(\sigma)\) be \(P\)'s result on \(\sigma\).

(\(\Leftarrow\))
\emph{Input slack.} Let \(\sigma_0 \notin \mathit{Pre}\) and let \(o\) be a
producible result different from \(o_P(\sigma_0)\). Patch \(P\) at
\(\sigma_0\). The patch agrees with \(P\) on \(\mathit{Pre}\), so it
satisfies the contract, and differs from \(P\) on \(\sigma_0\).

\emph{Output slack.} Let \(\sigma_0 \in \mathit{Pre}\) and
\(\theta = \theta_P\) with two accepted results, one of them \(o \neq
o_P(\sigma_0)\). (If the slack is only at another parameter, it is the
logical-parameter case.) Patch \(P\) at \(\sigma_0\) to produce \(o\).

\emph{Logical-parameter slack.} Let \(\sigma_0 \in \mathit{Pre}\) and
\(\theta_1, \theta_2 \in \mathcal{I}\) with
\(\mathit{Acc}_{\theta_1}(\sigma_0) \neq \mathit{Acc}_{\theta_2}(\sigma_0)\).
Both sets are nonempty. Some accepted result \(o\) at one of those
parameters differs from \(o_P(\sigma_0)\). Start from a program for that
parameter and patch it at \(\sigma_0\).

(\(\Rightarrow\)) If there is no such slack, then \(\mathit{Pre} = \Sigma\),
each \(\mathit{Acc}_\theta(\sigma)\) has at most one producible element, and
those elements coincide for all \(\theta \in \mathcal{I}\). That element is
\(o_P(\sigma)\). Any legal \(P'\) must produce it everywhere, so
\(P' \approx P\).
\end{proof}

Unexposed slack is omitted because it need not change observable API results.

\paragraph{Lemma~\ref{lem:lowering} (Lowering Preserves a Proved Contract).}
If \(P^* \models \Phi_{\text{func}}\) (respectively \(\models_K\)), then
\(\mathcal{L}_{\text{Target}}(P^*) \models \Phi_{\text{func}}\)
(respectively \(\models_K\)).
\begin{proof}
Let \(B = \mathcal{L}_{\text{Target}}(P^*)\). Parameters of \(B\) are those
of the source given to the backend. If \(P^*\) satisfies \(\Phi_{\text{func}}\)
(or \(\models_K\)) then it has no undefined behavior on the covered inputs
(or traces), terminates with an accepted result, and the compiler
correctness assumption copies that result to \(B\).
Proposition~\ref{prop:refine} is the unbounded case.
\end{proof}

\paragraph{Theorem~\ref{thm:restricted} (Soundness of a Restricted Intent Compiler).}
Let \(V\) be sound for \(\models\) or for \(\models_K\), and let the compiler
emit source \(P^*\) and binary
\(\mathrm{Bin}^* = \mathcal{L}_{\text{Target}}(P^*)\).
\begin{enumerate}
  \item If \(P^* \neq P_1\), then \(P^*\) is in \(\mathrm{Gen}(P_1, R, D)\),
  \(V\) proved it, and it passed \(\mathcal{B}\). Hence
  \(P^* \models \Phi_{\text{func}}\) or \(P^* \models_K \Phi_{\text{func}}\),
  matching \(V\), and the same judgment holds of \(\mathrm{Bin}^*\) by
  Lemma~\ref{lem:lowering}.
  \item The measured cost of \(\mathrm{Bin}^*\) is at most that of
  \(\mathcal{L}_{\text{Target}}(P_1)\), when measurement runs.
  \item If \(P_1 \models \Phi_{\text{func}}\) (or \(\models_K\)), then
  \(\mathrm{Bin}^*\) satisfies the same judgment in every case.
\end{enumerate}
\begin{proof}
Winner selection (Algorithm~\ref{alg:winner}) returns \(P_1\) or a member of \(\mathcal{A}\). If
\(P^* \neq P_1\), then \(P^* \in \mathcal{A}\), so \(P^* \in \mathrm{Gen}\),
\(V\) answered \(\mathsf{proved}\), and \(\mathcal{B}\) passed.
Soundness of \(V\) gives \(\models\) or \(\models_K\).
Lemma~\ref{lem:lowering} transfers the judgment to the binary. Measured
cost is minimized over a set containing
\(\mathcal{L}_{\text{Target}}(P_1)\). If \(P^* = P_1\) and \(P_1\) already
satisfies the judgment, apply Lemma~\ref{lem:lowering} to \(P_1\).
\end{proof}

\section{Pipeline Algorithms}
\label{app:algorithms}

This appendix gives each stage of \textsc{IntentCompile} as an algorithm:
inputs, outputs, the promises that algorithm makes, and a proof of those
promises. Every stage claims termination, relative to its oracles returning
(the solver for \(\mathcal{I}\), \(V\), the checker for \(\mathcal{B}\),
lowering, and measurement). Parameters such as the TheoryLower score, the
CEGIS distance, the attempt cap, the checker for \(\mathcal{B}\), and the
rewrite catalogs are inputs of the algorithms.
Section~\ref{sec:implementation} instantiates them. Hole domains are finite. \(V\) returns
\(\mathsf{proved}\), a counterexample, or \(\mathsf{unknown}\), and respects
its time budget. A rewrite in a catalog has a syntactic applicability
condition and optional side constraints on holes. Nothing below claims that
\(\mathrm{Sat}(\Phi_{\text{func}})\) is searchable, that the catalogs are
complete, or that \(V\) decides every query.

\subsection{\textsc{CheckRealiz}}

\begin{algo}[\textsc{CheckRealiz}]
\label{alg:checkrealiz}
Input: \(\Phi_{\text{func}} = \langle \mathit{Pre}, \mathit{Post},
\mathcal{I}\rangle\) and a solver for \(\mathcal{I}\) that answers
satisfiable, unsatisfiable, or \(\mathsf{unknown}\). A time budget on that
solver, when present, is mapped to \(\mathsf{unknown}\).
Output: \(\mathsf{continue}\), or a stop
(\(\mathsf{vacuous}\), \(\mathsf{unsat}\) with a core, or
\(\mathsf{unknown}\)).
\begin{enumerate}
  \item If \(\mathit{Pre} = \emptyset\), report \(\mathsf{vacuous}\) and
  stop.
  \item Ask the solver whether \(\mathcal{I}\) is satisfiable.
  \item If the answer is unsatisfiable, report \(\mathsf{unsat}\) with the
  solver's unsatisfiable core and stop.
  \item If the answer is \(\mathsf{unknown}\), stop.
  \item If the answer is satisfiable, report \(\mathsf{continue}\).
\end{enumerate}
\end{algo}

\begin{lemma}[\textsc{CheckRealiz}]
\label{lem:checkrealiz}
Assume the solver for \(\mathcal{I}\) is sound.
\begin{enumerate}
  \item \emph{Termination.} If the solver for \(\mathcal{I}\) returns in
  finite time, so does the algorithm. A timeout is \(\mathsf{unknown}\) and
  is a return.
  \item \emph{No candidate.} The algorithm emits no program.
  \item \emph{Vacuous stop.} If it reports \(\mathsf{vacuous}\), then
  \(\mathit{Pre} = \emptyset\), so every program satisfies
  \(\Phi_{\text{func}}\) and the contract states no intent.
  \item \emph{Unsat stop.} If it reports \(\mathsf{unsat}\), then
  \(\mathcal{I}\) is empty, so no program satisfies
  \(\Phi_{\text{func}}\).
  \item \emph{Unknown stop.} An \(\mathsf{unknown}\) answer is a stop, not a
  proof that \(\mathcal{I}\) is inhabited or empty.
  \item \emph{Continue.} A \(\mathsf{continue}\) result is only that
  \(\mathcal{I}\) has a model. It is not a candidate, and it does not decide
  whether \(\mathit{Post}\) is realizable on \(\mathit{Pre}\).
\end{enumerate}
\end{lemma}
\begin{proof}
Steps 1 and 3--5 are local. The only query is step 2. If the solver
returns satisfiable, unsatisfiable, or \(\mathsf{unknown}\), one of those
steps returns, which is (1). The algorithm has no step that constructs or
admits a program, so (2) holds. Step 1 returns \(\mathsf{vacuous}\) exactly
when \(\mathit{Pre} = \emptyset\). Satisfaction of \(\Phi_{\text{func}}\)
quantifies over \(\mathit{Pre}\), so an empty precondition is true of every
program, which is (3). Step 3 returns \(\mathsf{unsat}\) only on a sound
unsat answer, which means \(\mathcal{I}\) has no valuation. A program that
satisfies \(\Phi_{\text{func}}\) must have \(\theta_P \in \mathcal{I}\), so
no such program exists, which is (4). Step 4 returns on
\(\mathsf{unknown}\) without a model or a core, which is (5). The query in
step 2 mentions \(\mathcal{I}\) only. A model of \(\mathcal{I}\) does not
produce a program and does not prove that every \(\sigma \in \mathit{Pre}\)
has an accepted result under \(\mathit{Post}\), which is (6).
\end{proof}

The stronger query \(\forall \sigma \in \mathit{Pre}.\ \exists o.\
(\sigma, o) \in \mathit{Post}_\theta\) is not this stage. An unsatisfiable
harness is discovered later, when no candidate is proved. Termination is
relative to the solver: the stage does not decide an arbitrary theory of
\(\mathcal{I}\), and a solver with no budget may diverge.

\subsection{\textsc{TheoryLower}}

Let \(\mathit{score}\) map sketches to \(\delta\cdot\mathbb{N}\) for some
granularity \(\delta > 0\), and let \(\tau\) be a threshold at which
rewriting stops. Both are inputs. The intended reading is a weighted count
of constructs outside a chosen cheap set for the source language and
\(V\), so descent aims at that set. Any score with those values for which
every applied rewrite strictly decreases the score makes the lemma hold.
\textsc{IntC}'s score is the weighted feature count of
Section~\ref{sec:implementation}. \(\mathit{score}(S) \le \tau\) does not
imply that \(V\) decides \(S\).

\begin{algo}[\textsc{TheoryLower}]
\label{alg:theorylower}
Input: baseline \(P_1\), catalog \(R\), \(\mathit{score}\), and
threshold \(\tau\).
Output: a sketch \(S\) tagged \(\mathsf{below\text{-}threshold}\) or
\(\mathsf{stuck}\), or no sketch.
\begin{enumerate}
  \item Set \(S \gets P_1\). Rewrites, not this step, introduce holes.
  \item If \(\mathit{score}(S) \le \tau\), return no sketch, tagged
  \(\mathsf{below\text{-}threshold}\).
  \item Among pairs of a rewrite \(r \in R\) and a site of \(S\) where
  \(r\) applies, choose one that minimizes \(\mathit{score}(r(S))\). If no
  pair gives a strict decrease: if a rewrite has already been applied,
  return \((S, \mathsf{stuck})\); otherwise return no sketch, tagged
  \(\mathsf{stuck}\).
  \item Set \(S \gets r(S)\). If \(\mathit{score}(S) \le \tau\), return
  \((S, \mathsf{below\text{-}threshold})\). Otherwise go to step 3.
\end{enumerate}
\end{algo}

\begin{lemma}[\textsc{TheoryLower}]
\label{lem:theorylower}
\begin{enumerate}
  \item \emph{Termination.} The algorithm returns after at most
  \(\mathit{score}(P_1)/\delta\) rewrites.
  \item \emph{Catalog.} Every rewrite it applies is from \(R\), at a site
  that matches that rewrite. Holes come from those rewrites.
  \item \emph{Exits.} CEGIS receives a sketch only if a rewrite applied.
  That sketch has \(\mathit{score} \le \tau\), or no remaining rewrite in
  \(R\) lowers the score. If no rewrite applied, there is no TheoryLower
  seed.
  \item \emph{Unverified.} No returned sketch is submitted to \(V\) at this
  stage.
\end{enumerate}
\end{lemma}
\begin{proof}
Each applied rewrite lowers \(\mathit{score}\) by at least \(\delta\), and
the score is nonnegative, so that can happen at most
\(\mathit{score}(P_1)/\delta\) times, which is (1). Step 3 selects from
\(R\) at matching sites, and step 1 does not punch holes, which is (2).
Step 2 returns no sketch. Step 3 returns a sketch only if one was already
applied, or no sketch. Step 4 returns a rewritten sketch, which is (3).
No step calls \(V\), which is (4).
\end{proof}

The stage makes no semantic claim. The sketch may violate the contract for
every hole value, and later stages decide which fillings survive. The loop
is greedy: a rewrite that first raises the score and enables a larger drop
later is never taken. A \(\mathsf{stuck}\) result does not abort
\textsc{GeneratePool}, which starts from \(P_1\).

\subsection{\textsc{Baseline-Directed CEGIS}}

Let \(S[h]\) be a sketch with holes ranging over a finite domain \(D\), and
let \(\mathit{dist}\) be a distance on \(D\). Closeness is relative to
\(\mathit{dist}\). \textsc{IntC} minimizes distance of baseline holes to
\(\theta_{P_1}\) and extra holes subject to the side constraints
(Section~\ref{sec:implementation}).

\begin{algo}[\textsc{Baseline-Directed CEGIS}]
\label{alg:cegis}
Input: sketch \(S\), domain \(D\), side constraints, \(\mathcal{I}\),
baseline parameters \(\theta_{P_1}\), verifier \(V\) with budget \(t\), and
a finite attempt cap \(N\).
Output: a filled program \(S[h]\), or \(\mathsf{fail}\).
\begin{enumerate}
  \item Set \(R_{\mathrm{ex}} \gets \emptyset\) and \(n \gets 0\).
  \item If \(n = N\), return \(\mathsf{fail}\).
  \item Find \(h \in D \setminus R_{\mathrm{ex}}\) that meets the side
  constraints, gives \(\theta_{S[h]} \in \mathcal{I}\), and minimizes
  \(\mathit{dist}(h, \theta_{P_1})\). If none exists, return
  \(\mathsf{fail}\).
  \item Run \(V(S[h], \Phi_{\text{func}}, t)\).
  \item If the answer is \(\mathsf{proved}\), return \(S[h]\).
  \item If the answer is a counterexample, add \(h\) to
  \(R_{\mathrm{ex}}\). Do not add the failing input as a constraint on
  later proposals. Set \(n \gets n+1\) and go to step 2.
  \item If the answer is \(\mathsf{unknown}\), return \(\mathsf{fail}\).
\end{enumerate}
\end{algo}

\paragraph{Theorem~\ref{thm:cegis} (Properties of Baseline-Directed CEGIS).}
Let \(S[h]\) range over a finite hole domain \(D\), and let \(N\) be the
attempt cap (\(N = |D|\) if there is no cap).
\begin{enumerate}
  \item If CEGIS returns \(S[h]\), then \(V\) proved \(S[h]\).
  \item \emph{Termination.} If \(V\) returns on each query, CEGIS returns
  after at most \(\min(N, |D|)\) iterations. Step 3 searches a finite set,
  so it returns.
  \item If some \(h^\circ \in D\) makes \(S[h^\circ]\) satisfy the contract
  (or \(\models_K\)), \(V\) never answers \(\mathsf{unknown}\), and there is
  no attempt cap, then CEGIS returns a satisfying assignment closest to
  \(\theta_{P_1}\).
\end{enumerate}
\begin{proof}
(1) The only success return is step 5, which requires
\(\mathsf{proved}\).

(2) \(V\) is assumed to return (\(\mathsf{proved}\), a counterexample, or
\(\mathsf{unknown}\)) because it respects \(t\). Each unsuccessful
iteration that does not fail on \(\mathsf{unknown}\) adds one \(h\) to
\(R_{\mathrm{ex}}\) and never retries it. The loop also stops when
\(n = N\). Together that is at most \(\min(N, |D|)\) iterations. Step 3
enumerates a finite remainder of \(D\).

(3) Let \(H^\circ = \{h' \in D \mid S[h'] \models \Phi_{\text{func}}\}\)
(or \(\models_K\)), nonempty by assumption. Soundness of \(V\) keeps every
member of \(H^\circ\) out of \(R_{\mathrm{ex}}\), because only a
counterexample excludes a valuation. Step 3 therefore always sees a
nonempty feasible set, so it does not fail. Since \(V\) never answers
\(\mathsf{unknown}\) and there is no cap, (2) ends the run on
\(\mathsf{proved}\). The returned \(h\) is feasible for step 3 at that
iteration, hence minimizes \(\mathit{dist}(\cdot, \theta_{P_1})\) over a
set that still contains \(H^\circ\). An attempt cap preserves (1) and (2)
and may drop (3).
\end{proof}

The bound \(|D|\) is a worst case. Side constraints shrink the search.
The guarantee covers only the holes of the given sketch. If no filling of
that sketch satisfies the contract, CEGIS fails, and a different structure
is a rewrite, not a closer \(h\).

\subsection{\textsc{BudgetNormalize}}

\begin{algo}[\textsc{BudgetNormalize}]
\label{alg:budgetnormalize}
Input: a finite set of CEGIS seeds, baseline \(P_1\), and a terminating
checker for \(\mathcal{B}\).
Output: a (possibly empty) subset of the seeds, and \(P_1\) unmarked as a
seed. \textsc{GeneratePool} is not skipped.
\begin{enumerate}
  \item Apply the checker for \(\mathcal{B}\) to each seed and to
  \(P_1\).
  \item Keep those seeds for which the checker passes. Drop the rest.
  \item If none remain, record that the seed set is empty.
  \item Return the kept set. Do not disable \textsc{GeneratePool}.
\end{enumerate}
\end{algo}

\begin{lemma}[\textsc{BudgetNormalize}]
\label{lem:budgetnormalize}
\begin{enumerate}
  \item \emph{Termination.} If the checker terminates on each of finitely
  many programs, so does the algorithm.
  \item \emph{Filter.} It adds no program.
  \item \emph{Kept seeds.} Every kept seed still carries whatever \(V\)
  proof it had, and passes \(\mathcal{B}\).
  \item \emph{Baseline.} The checker is applied to \(P_1\), but \(P_1\) is
  never added to the kept seeds, so a miss does not remove it from
  \eqref{eq:winner}.
  \item \emph{Empty seeds.} An empty kept set does not skip
  \textsc{GeneratePool}.
\end{enumerate}
\end{lemma}
\begin{proof}
One checker run per seed and one for \(P_1\) is finite if each run is,
which is (1). The algorithm only retains or drops members of its input,
which is (2). A kept seed is a subset of the input, so a prior
\(\mathsf{proved}\) answer remains, and step 2 is exactly the checker for
\(\mathcal{B}\), which is (3). Step 1 runs the checker on \(P_1\) and
step 2 never adds it to the kept set, which is (4). Step 4 returns after recording
emptiness and does not disable the next stage, which is (5).
\end{proof}

The stage does not repair a seed that misses \(\mathcal{B}\). A later
rewrite in \textsc{GeneratePool} may still produce a different program that
fits.

\subsection{\textsc{GeneratePool}}

Let \(R_{\mathrm{pool}}\) be a finite catalog of representation rewrites,
disjoint in purpose from the TheoryLower catalog though both are subsets of
the generator's \(R\).

\begin{algo}[\textsc{GeneratePool}]
\label{alg:generatepool}
Input: baseline \(P_1\), catalog \(R_{\mathrm{pool}}\), any seeds that
survived \textsc{BudgetNormalize}, verifier \(V\), and the checker for
\(\mathcal{B}\).
Output: a finite set \(\mathcal{C}\) of programs.
\begin{enumerate}
  \item Set \(\mathcal{C}\) to the surviving seeds (empty if none).
  \item For each rewrite \(r \in R_{\mathrm{pool}}\) that applies to
  \(P_1\), form the sketch \(r(P_1)\), run
  Algorithm~\ref{alg:cegis} on it, and if CEGIS returns a program that
  \(V\) proved and that passes \(\mathcal{B}\), add that program to
  \(\mathcal{C}\).
  \item Return \(\mathcal{C}\). Do not rank by \(\mathcal{O}\) or by an
  estimated cost.
\end{enumerate}
\end{algo}

\begin{lemma}[\textsc{GeneratePool}]
\label{lem:generatepool}
\begin{enumerate}
  \item \emph{Termination.} The algorithm terminates if each CEGIS run
  does and the checker for \(\mathcal{B}\) returns on each tried program.
  \item \emph{Starts from \(P_1\).} Step 2 runs even when the seed set is
  empty, and even when TheoryLower returned \(\mathsf{stuck}\).
  \item \emph{Admission.} Every program in the returned \(\mathcal{C}\)
  is in \(\mathcal{A}\). Surviving seeds are already in \(\mathcal{A}\).
  \item \emph{No ranking.} The algorithm does not consult
  \(\mathcal{O}\).
\end{enumerate}
\end{lemma}
\begin{proof}
\(R_{\mathrm{pool}}\) is finite. Each CEGIS run terminates by
Theorem~\ref{thm:cegis}(2). Each budget check is one run of a terminating
checker, which is (1). Step 2
iterates that catalog from \(P_1\) with no guard on the seed set or on
TheoryLower's tag, which is (2). A program added in step 2 was obtained
from \(P_1\) by a rewrite from the catalog, filled by CEGIS, proved by
\(V\), and passed \(\mathcal{B}\), hence lies in \(\mathcal{A}\). Surviving
seeds were proved by an earlier CEGIS run and passed
\textsc{BudgetNormalize}, which is (3). Step 3 forbids ranking, and no
other step reads \(\mathcal{O}\), which is (4).
\end{proof}

The forward pass can miss a combination that pays only after a later
rewrite. Exploring sequences of rewrites would find those combinations at
a cost exponential in sequence length. Ranking remains
\eqref{eq:winner}.

\subsection{Winner Selection}

\begin{algo}[Winner Selection]
\label{alg:winner}
Input: \(P_1\), the set \(\mathcal{C}\) from
Algorithm~\ref{alg:generatepool}, backend
\(\mathcal{L}_{\text{Target}}\), and objective \(\mathcal{O}\).
Output: source \(P^*\) and binary
\(\mathrm{Bin}^* = \mathcal{L}_{\text{Target}}(P^*)\).
\begin{enumerate}
  \item Let \(\mathcal{A}\) be \(\mathcal{C}\) (every member already
  proved and inside \(\mathcal{B}\)).
  \item If the contract names no \(\mathcal{O}\), return the first member
  of \(\mathcal{A}\), or \(P_1\) if \(\mathcal{A}\) is empty.
  \item Lower \(P_1\) and every member of \(\mathcal{A}\) with
  \(\mathcal{L}_{\text{Target}}\).
  \item Measure each binary on the named workload and hardware, and
  return the source whose binary minimizes \(\mathcal{O}\), preferring
  \(P_1\) on ties, as in~\eqref{eq:winner}. If \(\mathcal{A}\) is empty,
  return \(P_1\).
\end{enumerate}
\end{algo}

\begin{lemma}[Winner Selection]
\label{lem:winner}
\begin{enumerate}
  \item \emph{Termination.} If lowering and measurement return on each of
  finitely many programs, so does the algorithm. If there is no
  \(\mathcal{O}\), there is no measurement loop.
  \item It returns \(P_1\) or a member of \(\mathcal{A}\).
  \item Measurement is not admission: a program not proved by \(V\) is not
  ranked. With no \(\mathcal{O}\), the first admitted program is emitted.
\end{enumerate}
The remaining claims are Theorem~\ref{thm:restricted}.
\end{lemma}
\begin{proof}
\(\mathcal{A}\) is finite by Lemma~\ref{lem:generatepool}. With no
\(\mathcal{O}\), step 2 returns immediately, which is (1). Otherwise one
lowering and one measurement per member of
\(\mathcal{A} \cup \{P_1\}\) is finite if each call returns, which is (1).
The \(\arg\min\) set is
\(\{\mathcal{L}_{\text{Target}}(P) \mid P \in \mathcal{A} \cup \{P_1\}\}\),
or \(\{P_1\}\) when \(\mathcal{A}\) is empty, or the first of
\(\mathcal{A}\) when there is no \(\mathcal{O}\), so the returned source is
\(P_1\) or a member of \(\mathcal{A}\), which is (2). Steps 2 and 4 rank
only admitted programs, which is (3).
\end{proof}
